\chapter{Năng lượng chân không và chân trời de Sitter}

\lectureoverview{Năng lượng chân không khác với vật chất thông thường: nó không chọn ra một hệ quy chiếu quán tính ưu tiên nào và cũng không bị pha loãng khi vũ trụ giãn nở. Một khi có tính đến hấp dẫn, mốc không của năng lượng tưởng như vô hại ấy sẽ làm thay đổi hình học của không-thời gian. Chúng ta sẽ suy ra định luật chi phối bởi vật chất \(a(t)\propto t^{2/3}\), đối chiếu nó với định luật chi phối bởi chân không \(a(t)\propto e^{Ht}\), rồi sau đó viết lại không-thời gian de Sitter trong tọa độ bảo giác. Biểu đồ nhân quả thu được cho thấy một chân trời sự kiện phụ thuộc vào người quan sát ở khoảng cách riêng cố định.}

\section{Thăng giáng lượng tử trong chân không}

Chúng ta bắt đầu bằng câu hỏi thôi thúc toàn bộ bài giảng. Lý thuyết trường lượng tử gán năng lượng cho không gian trống rỗng. Điều đó có nghĩa là các cặp electron--positron thực liên tục xuất hiện từ hư vô hay không?

\begin{classroomqa}
\audiencequestion{Liệu sự tạo thành và hủy cặp electron--positron có thực sự xảy ra trong chân không, hay đó chỉ đơn thuần là một thủ pháp hạch toán?\lecturetimestamp{00:00:06}}
\lecturerresponse{Các cặp ảo góp phần vào năng lượng chân không, nhưng chúng không phải là những hạt thông thường mà ta có thể quan sát khi chúng xuất hiện từ không gian trống rỗng. Chúng là các thăng giáng của trạng thái cơ bản lượng tử và chỉ tồn tại trong những giới hạn thích hợp với một quá trình ảo.}
\end{classroomqa}

Dao động tử điều hòa cho ta mô hình minh bạch nhất. Về mặt cổ điển, nó có thể đứng yên tại điểm cực tiểu của thế năng. Theo cơ học lượng tử, vị trí chính xác tại điểm cực tiểu và động lượng đúng bằng 0 là không tương thích. Trạng thái cơ bản dung hòa giữa hai điều đó và vẫn giữ năng lượng điểm không
\begin{equation}
E_0=\frac{1}{2}\hbar\omega.
\label{eq:l08-zero-point}
\end{equation}
Đây là năng lượng nhỏ nhất mà dao động tử có thể có. Nó không phải là trạng thái kích thích và không biểu diễn một dòng hạt có thể quan sát được.

Một trường tự do phân tích thành các mode chuẩn, mỗi mode hành xử như một dao động tử. Trường điện từ trong một hốc cộng hưởng là một ví dụ quen thuộc. Mỗi mode đóng góp một hạng có dạng như trong Phương trình~\eqref{eq:l08-zero-point}; các trường lượng tử của electron, quark và các trường khác đóng góp các thăng giáng trạng thái cơ bản riêng của chúng. Cộng các đóng góp này lại tạo ra cái mà ta gọi là năng lượng chân không. Sự không khớp tai tiếng giữa các ước lượng vi mô ngây thơ và thang vũ trụ học được quan sát là có thật, nhưng đó không phải là chủ đề của bài giảng này.

\begin{classroomqa}
\audiencequestion{Liệu sự sai khác cỡ \(120\) bậc độ lớn có được thảo luận lại không?\lecturetimestamp{00:02:59}}
\lecturerresponse{Không ở đây. Mục tiêu hiện tại là hiểu các chân trời vũ trụ: các phương trình vũ trụ học tạo ra sự giãn nở gia tốc như thế nào và vì sao sự giãn nở gia tốc lại tạo ra các chân trời sự kiện.}
\end{classroomqa}

\begin{classroomqa}
\audiencequestion{Quan hệ tử ngoại--hồng ngoại từ bài giảng trước có phải là một phần của thảo luận tối nay không?\lecturetimestamp{00:03:59}}
\lecturerresponse{Có thể sẽ quay lại sau, nhưng nó không nằm trong mạch lập luận đã định. Lộ trình trước mắt đi từ năng lượng chân không đến vũ trụ học rồi đến các chân trời.}
\end{classroomqa}

\section{Vì sao năng lượng chân không là bất biến Lorentz}

Câu hỏi tiếp theo có vẻ sơ đẳng nhưng phơi bày tính chất định nghĩa của năng lượng chân không. Nếu không gian trống rỗng có mật độ năng lượng dương, tại sao nó không hành xử như một môi trường quang học và làm thay đổi tốc độ ánh sáng?

\begin{classroomqa}
\audiencequestion{Liệu một mật độ năng lượng chân không dương có nên có hiệu ứng khúc xạ, khiến ánh sáng truyền đi chậm hơn đôi chút không?\lecturetimestamp{00:04:35}}
\lecturerresponse{Vật chất thông thường có thể làm như vậy vì nó xác định một hệ quy chiếu nghỉ ưu tiên. Năng lượng chân không thì khác: nó không chọn ra một hệ như thế, và mọi quan sát viên quán tính đều thu được cùng một mật độ năng lượng chân không. Vì vậy, các phép đo cục bộ của ánh sáng vẫn giữ tốc độ bất biến \(c\).}
\end{classroomqa}

Ta phải phân biệt đối xứng của các định luật với đối xứng của một trạng thái cụ thể. Các phương trình vật lý có thể bất biến tịnh tiến mặc dù Stanford không phải Berkeley và Trái Đất không phải không gian liên sao. Tịnh tiến \emph{toàn bộ} sự sắp đặt vật lý sẽ để các định luật không đổi; còn sự sắp đặt vật chất thực tế thì không nhất thiết phải trông giống nhau ở mọi vị trí. Tương tự, các phương trình bất biến quay không làm cho phòng học này có đối xứng quay.

Sự phân biệt tương tự cũng áp dụng cho các phép biến đổi Lorentz. Một chất khí có một hệ quy chiếu nghỉ. Một quan sát viên chuyển động xuyên qua nó sẽ thấy các phân tử chuyển động và một mật độ năng lượng khác, nên trạng thái của chất khí không bất biến Lorentz mặc dù các phương trình chi phối nó thì có. Vật chất thông thường, các môi trường bức xạ và các thiết bị trong phòng thí nghiệm đều phá vỡ đối xứng bằng cách chọn ra những chuyển động hoặc phương đặc biệt.

Năng lượng chân không là dạng vật chất ngoại lệ không làm như thế. Hãy tưởng tượng một đầu dò cho biết mật độ năng lượng chân không. Khi đứng yên, nó cho một giá trị nào đó \(\rho_{\rm vac}\). Gia tốc cho đầu dò rồi lặp lại phép đo: kết quả không đổi. Trong ký hiệu tương đối tính, tính chất này được mã hóa bằng một tensor năng lượng-xung lượng tỉ lệ với metric,
\begin{equation}
T_{\mu\nu}^{\rm vac}=-\rho_{\rm vac}g_{\mu\nu},
\qquad
p_{\rm vac}=-\rho_{\rm vac}.
\label{eq:l08-vacuum-stress}
\end{equation}
\editorialnote{Quan hệ áp suất không được suy ra trên bảng. Đó là phát biểu hiệp biến đứng sau hai tính chất định nghĩa của bài giảng: bất biến Lorentz và mật độ không đổi trong quá trình giãn nở.}

Điều kiện cục bộ là quan trọng. Một tia sáng đi qua một quan sát viên được đo bằng các đồng hồ và thước đo lân cận là chuyển động với \(c\). Việc hỏi về tốc độ của một tia sáng ở xa đối với ta trong một hình học giãn nở đòi hỏi phải chọn tọa độ và một quy ước để so sánh các sự kiện ở xa. Tốc độ theo tọa độ đó không nhất thiết bằng \(c\), mặc dù mọi quan sát viên cục bộ đều đo được \(c\).

\begin{classroomqa}
\audiencequestion{Trong thảo luận này, năng lượng chân không và hằng số vũ trụ có đồng nghĩa không?\lecturetimestamp{00:11:57}}
\lecturerresponse{Có. Chúng là hai cách đặt cùng một hạng tử trong phương trình Einstein: nó có thể được mô tả như một hằng số vũ trụ ở phía hình học hoặc như năng lượng-xung lượng chân không bất biến Lorentz ở phía nguồn.}
\end{classroomqa}

\section{Cái gì bị pha loãng và cái gì không bị pha loãng}

Hãy lấp đầy một vùng lớn bằng các hạt không tương đối tính có khối lượng \(m\). Trên các thang đủ lớn để lấy trung bình qua các thiên hà và cụm thiên hà, giả sử phân bố là đồng nhất. Nếu một thể tích đồng chuyển động chứa \(N\) hạt, thì mật độ vật chất của nó là
\begin{equation}
\rho_m=\frac{Nm}{V}.
\end{equation}
Khi mọi kích thước tuyến tính tăng gấp đôi, thể tích tăng \(2^3=8\), còn \(N\) và \(m\) vẫn không đổi. Vì vậy mật độ giảm xuống còn một phần tám giá trị trước đó. Vật chất thông thường bị pha loãng vì cùng những hạt ấy chiếm một thể tích vật lý lớn hơn.

Bức xạ cũng bị pha loãng, và còn nhanh hơn. Mật độ số của nó giảm theo thể tích, trong khi bước sóng của mỗi photon bị kéo giãn và năng lượng của nó bị dịch đỏ. Chúng ta sẽ ghi lại kết quả \(\rho_r\propto a^{-4}\) sau, nhưng tối nay sẽ không giải vũ trụ học chi phối bởi bức xạ.

Năng lượng chân không là ngoại lệ:
\begin{equation}
\rho_{\rm vac}(t)=\text{constant}.
\label{eq:l08-rho-vac-constant}
\end{equation}
Nó không kết tụ, không chọn ra một hướng, không chọn ra một hệ quy chiếu quán tính, và không giảm mật độ khi không gian giãn nở. Nếu một hộp đồng chuyển động lớn lên, năng lượng chân không chứa trong đó tăng tỉ lệ với thể tích của hộp.

\editorialnote{Mật độ không đổi không có nghĩa là tổng năng lượng chân không bên trong một hộp đồng chuyển động đang giãn nở là không đổi. Bảo toàn hiệp biến năng lượng-xung lượng,
\(\dot\rho+3H(\rho+p)=0\), được thỏa mãn vì năng lượng chân không có áp suất âm \(p_{\rm vac}=-\rho_{\rm vac}\). Hạng công cung cấp chính xác độ biến thiên của tổng năng lượng chân không trong hộp.}

\section{Hấp dẫn nhìn thấy mốc không của năng lượng}

Trong vật lý phòng thí nghiệm không có hấp dẫn, việc dịch chuyển mọi năng lượng cùng một hằng số thường không làm thay đổi điều gì. Phổ phụ thuộc vào các hiệu số: năng lượng giữa một nguyên tử ở trạng thái kích thích và không kích thích, hoặc giữa một trạng thái có hạt và một trạng thái không có hạt. Không gian trống rỗng dường như không cho một mốc tham chiếu nào để đo năng lượng tuyệt đối của nó.

Hấp dẫn làm thay đổi câu hỏi. Khối lượng là năng lượng, năng lượng là nguồn của hấp dẫn, và vì vậy năng lượng chân không cũng hấp dẫn. Một hằng số vốn vô hình trong một chuyển mức nguyên tử có thể làm cong không-thời gian.

Phương trình Einstein tách hình học khỏi các nguồn:
\begin{equation}
G_{\mu\nu}=8\pi G_N T_{\mu\nu},
\label{eq:l08-einstein}
\end{equation}
trong các đơn vị với \(c=1\). Tensor Einstein \(G_{\mu\nu}\) được xây dựng từ metric \(g_{\mu\nu}(x)\) và các đạo hàm của nó. Tensor năng lượng-xung lượng \(T_{\mu\nu}\) chứa mật độ năng lượng, mật độ xung lượng, áp suất và ứng suất do vật chất và các trường không hấp dẫn gây ra. Năng lượng chân không thuộc về phía nguồn này cũng như vật chất thông thường.\editorialnote{Phần dẫn nhập bằng lời chỉ liên hệ ngắn gọn \(8\pi/3\) với phương trình tensor. Thừa số \(1/3\) thuộc về phương trình Friedmann phẳng không gian ở dưới, không phải Phương trình~\eqref{eq:l08-einstein}.}

\begin{figure}[htbp]
\centering
\includegraphics[width=0.62\textwidth]{lecture_08_figure_02.png}
\caption{Vế hình học của phương trình Einstein, được xây dựng từ metric không-thời gian \(g_{\mu\nu}(x)\).}
\lectureframe{00:21:58.820}
\end{figure}

\section{Hình học quy về một thừa số tỉ lệ}

Bây giờ ta chọn hình học giãn nở đồng nhất đơn giản nhất với các lát không gian phẳng:
\begin{equation}
ds^2=-dt^2+a(t)^2\,d\mathbf{x}^2,
\qquad
d\mathbf{x}^2=dx^2+dy^2+dz^2.
\label{eq:l08-frw-metric}
\end{equation}
Các mốc tọa độ \(\mathbf{x}\) được vẽ trên một tấm cao su vô hạn. Bản thân chúng không nhất thiết chuyển động, nhưng khoảng cách vật lý giữa chúng thì có thể thay đổi. Hai điểm đồng chuyển động cách nhau \(\Delta\mathbf{x}\) có khoảng cách riêng tại \(t\) cố định là
\begin{equation}
D(t)=a(t)\lvert\Delta\mathbf{x}\rvert.
\label{eq:l08-proper-distance}
\end{equation}
Vì vậy, toàn bộ phụ thuộc theo thời gian của hình học đối xứng cao này được chứa trong một hàm duy nhất, thừa số tỉ lệ \(a(t)\). \(a\) không đổi cho không-thời gian Minkowski phẳng sau một phép co giãn không gian. \(a\) phụ thuộc thời gian có thể giữ cho mỗi lát không gian phẳng trong khi vẫn làm cho toàn bộ không-thời gian cong.

\begin{classroomqa}
\audiencequestion{Với các mốc tọa độ \(x=7\) và \(x=13\), khoảng cách tọa độ của chúng có phải là \(4\) không?\lecturetimestamp{00:27:40}}
\lecturerresponse{Không. Nó là \(13-7=6\), nên khoảng cách vật lý của chúng là \(6a(t)\). Việc sửa lại này rất hữu ích: \(x\) gắn nhãn vị trí đồng chuyển động, còn nhân với \(a(t)\) thì cho khoảng cách vật lý.}
\end{classroomqa}

Đối với metric này và một mật độ năng lượng đồng đều theo không gian, một thành phần của phương trình Einstein trở thành phương trình Friedmann phẳng không gian,
\begin{equation}
\left(\frac{\dot a}{a}\right)^2
=\frac{8\pi G_N}{3}\rho(t).
\label{eq:l08-friedmann}
\end{equation}
Tỉ số
\begin{equation}
H(t)\equiv\frac{\dot a(t)}{a(t)}
\label{eq:l08-hubble-definition}
\end{equation}
là tham số Hubble. Ở giai đoạn này, đó là một định nghĩa và có thể biến thiên theo thời gian.

\begin{figure}[htbp]
\centering
\includegraphics[width=0.90\textwidth]{lecture_08_figure_03.png}
\caption{Metric giãn nở phẳng và phương trình Friedmann sau khi thừa số còn thiếu \(1/a^2\) đã được khôi phục trên bảng.}
\lectureframe{00:30:30.000}
\end{figure}

\paragraph{Bức tranh cân bằng năng lượng.}
Đối với bức tranh cân bằng năng lượng này, Phương trình~\eqref{eq:l08-friedmann} được viết một cách sơ đồ thành
\begin{equation}
-\frac{3}{8\pi G_N}H^2+\rho=0.
\label{eq:l08-energy-heuristic}
\end{equation}
Năng lượng vật chất là dương, còn bậc tự do hấp dẫn được biểu diễn bởi thừa số tỉ lệ đang giãn nở mang một hạng động năng âm. Trong bài toán rút gọn này, việc năng lượng hấp dẫn âm cộng với năng lượng vật chất dương bằng 0 là một mẹo ghi nhớ hữu ích.

\editorialnote{Phương trình~\eqref{eq:l08-energy-heuristic} là ràng buộc Hamilton trong một mô hình minisuperspace có đối xứng rất cao, chứ không phải một định lý độc lập với tọa độ rằng mọi vũ trụ bất kỳ đều có một tổng năng lượng toàn cục được bảo toàn bằng 0. Bài giảng dùng nó như một cách diễn giải trực giác của ràng buộc Friedmann.}

\section{Sự giãn nở do vật chất chi phối}

Hãy quay lại một hộp toán học đồng chuyển động có cạnh tọa độ \(\Delta x=1\). Cạnh vật lý của nó là \(a(t)\), nên thể tích vật lý của nó là \(a(t)^3\). Nếu nó chứa \(N\) hạt không tương đối tính có khối lượng \(m\), thì
\begin{equation}
\rho_m(t)=\frac{Nm}{a(t)^3}
=\frac{\rho_0}{a(t)^3}.
\label{eq:l08-matter-density}
\end{equation}
Bản thân các hạt không giãn nở; khoảng cách trung bình giữa chúng thì có.

\begin{figure}[htbp]
\centering
\includegraphics[width=0.84\textwidth]{lecture_08_figure_04.png}
\caption{Metric và phương trình Friedmann khi lập luận về hộp đồng chuyển động được thiết lập.}
\lectureframe{00:36:07.940}
\end{figure}

Thế vào Phương trình~\eqref{eq:l08-friedmann} ta được
\begin{equation}
\left(\frac{\dot a}{a}\right)^2
=\frac{8\pi G_N\rho_0}{3a^3}.
\end{equation}
Đặt \(K=8\pi G_N\rho_0/3\) và chọn nhánh giãn nở:
\begin{align}
\dot a^2&=\frac{K}{a},\\
\dot a&=\sqrt{K}\,a^{-1/2},\\
a^{1/2}\,da&=\sqrt{K}\,dt.
\label{eq:l08-matter-separated}
\end{align}

\begin{classroomqa}
\audiencequestion{Phương trình đã tách biến này nên được giải như thế nào?\lecturetimestamp{00:38:15}}
\lecturerresponse{Lấy tích phân hai vế. Tích phân của \(a^{1/2}\,da\) là \(\frac{2}{3}a^{3/2}\), không phải \(\frac{3}{2}a^{3/2}\); \(\frac{3}{2}\) là số mũ mới, còn \(\frac{2}{3}\) là hệ số.}
\end{classroomqa}

Do đó,
\begin{equation}
\frac{2}{3}a^{3/2}
=\sqrt{K}\,(t-t_B),
\end{equation}
trong đó \(t_B\) là một hằng số tích phân ấn định gốc thời gian đã chọn. Ngoài chuẩn hóa và độ dời thời gian,
\begin{equation}
a(t)\propto (t-t_B)^{2/3}.
\label{eq:l08-matter-solution}
\end{equation}

\begin{figure}[htbp]
\centering
\includegraphics[width=0.86\textwidth]{lecture_08_figure_05.png}
\caption{Phép thế mật độ vật chất và phép tách biến dẫn đến \(a^{1/2}da\propto dt\).}
\lectureframe{00:38:30.000}
\end{figure}

Đường cong \(t^{2/3}\) tăng chậm hơn một đường thẳng. Độ dốc của nó giảm dần, nên một vũ trụ bị vật chất chi phối là vũ trụ giảm tốc. Theo ngôn ngữ Newton được dùng trong bài giảng, các hạt ban đầu tách xa nhau, và lực hấp dẫn lẫn nhau của chúng liên tục làm giảm tốc độ tách xa.

\begin{figure}[htbp]
\centering
\includegraphics[width=0.86\textwidth]{lecture_08_figure_06.png}
\caption{So sánh trên bảng giữa sự giãn nở tuyến tính và quy luật vật chất lõm xuống \(a(t)\propto t^{2/3}\).}
\lectureframe{00:40:00.000}
\end{figure}

Để định hướng, một phông bức xạ có
\begin{equation}
\rho_r\propto a^{-4},
\end{equation}
chứ không phải \(a^{-3}\): sự pha loãng do thể tích đóng góp ba lũy thừa và độ dịch đỏ vũ trụ học đóng góp thêm một lũy thừa nữa. Hệ số tỉ lệ tương ứng cũng có thể giải được, nhưng bài giảng chủ ý để riêng bài tập đó.

\section{Sự giãn nở bị năng lượng chân không chi phối}

Bây giờ chỉ đặt năng lượng chân không ở vế phải của phương trình Friedmann. Vì \(\rho_{\rm vac}\) là hằng số,
\begin{equation}
\left(\frac{\dot a}{a}\right)^2
=\frac{8\pi G_N}{3}\rho_{\rm vac}
\equiv H^2.
\label{eq:l08-de-sitter-friedmann}
\end{equation}
Cần phân biệt hai cách dùng của \(H\). Phương trình~\eqref{eq:l08-hubble-definition} định nghĩa đại lượng có thể phụ thuộc thời gian \(\dot a/a\). Phương trình~\eqref{eq:l08-de-sitter-friedmann} nói rằng trong một vũ trụ thuần năng lượng chân không, giá trị của nó là hằng số
\begin{equation}
H=\sqrt{\frac{8\pi G_N\rho_{\rm vac}}{3}}.
\end{equation}
Chọn nhánh giãn nở thì được
\begin{equation}
\dot a=Ha,
\end{equation}
và do đó
\begin{equation}
a(t)=C e^{Ht}.
\label{eq:l08-exponential-solution}
\end{equation}
Hằng số tích phân \(C\) xác định sự chuẩn hóa tùy ý của các tọa độ đồng chuyển động. Việc tăng lên tỉ lệ với giá trị hiện thời của \(a\) chính là tăng trưởng mũ.

\begin{figure}[htbp]
\centering
\includegraphics[width=0.88\textwidth]{lecture_08_figure_07.png}
\caption{Phương trình bị năng lượng chân không chi phối \((\dot a/a)^2=H^2\), dạng bậc nhất của nó \(\dot a=Ha\), và nghiệm mũ.}
\lectureframe{00:47:45.000}
\end{figure}

Định luật Hubble cho các vật thể đồng chuyển động là
\begin{equation}
v_{\rm rec}=HD.
\label{eq:l08-hubble-law}
\end{equation}
Trong trường hợp de Sitter thuần túy này, khoảng cách riêng tại đó tốc độ lùi bằng \(c\) là hằng số:
\begin{equation}
D_H=\frac{c}{H},
\label{eq:l08-hubble-radius}
\end{equation}
hay \(D_H=H^{-1}\) trong các đơn vị \(c=1\) dùng trên bảng. Thang độ cố định này sẽ lại xuất hiện về mặt hình học như bán kính chân trời sự kiện de Sitter.\editorialnote{Bán kính Hubble không tự động là chân trời sự kiện trong một vũ trụ học phụ thuộc thời gian bất kỳ. Sự đồng nhất này đúng ở đây vì \(H\) là hằng số trong không gian de Sitter chính xác.}

Bài giảng nêu ra một thí nghiệm tưởng tượng cho máy dò năng lượng chân không. Hãy đặt hai vật thử rất nhẹ ở xa các nguồn khác, ban đầu đứng yên đối với dòng đồng chuyển động, và làm cho lực hút lẫn nhau của chúng không đáng kể. Khoảng cách giữa chúng tuân theo sự giãn nở mũ. Hiệu ứng này quá nhỏ đối với một thiết bị trên Trái Đất, nhưng các khoảng cách vũ trụ và thời gian nhìn ngược làm cho cùng vật lý ấy có thể quan sát được.

\section{Từ sự chi phối của vật chất sang sự chi phối của năng lượng chân không}

Mô hình vũ trụ học thực tế được xét trong bài giảng chứa cả hai hạng:
\begin{equation}
H(t)^2
=\frac{8\pi G_N}{3}
\left(\frac{\rho_0}{a(t)^3}+\rho_{\rm vac}\right).
\label{eq:l08-mixed-friedmann}
\end{equation}
Khi \(a\) nhỏ, hạng vật chất chi phối vì \(a^{-3}\) lớn, và nghiệm giống với \(t^{2/3}\). Khi vũ trụ lớn lên, vật chất bị pha loãng còn năng lượng chân không vẫn giữ nguyên. Cuối cùng hạng hằng chiếm ưu thế và hệ số tỉ lệ uốn cong đi lên hướng tới \(e^{Ht}\).

\begin{figure}[htbp]
\centering
\includegraphics[width=0.87\textwidth]{lecture_08_figure_08.png}
\caption{Chuyển tiếp định tính được vẽ trong bài giảng: giai đoạn đầu giảm tốc do vật chất chi phối, tiếp theo là tăng trưởng mũ muộn do năng lượng chân không chi phối.}
\lectureframe{00:53:30.000}
\end{figure}

\begin{classroomqa}
\audiencequestion{Nghiệm mũ chứa một hằng số tích phân. Làm sao biết rằng nó không cực kỳ nhỏ?\lecturetimestamp{00:52:41}}
\lecturerresponse{Ta đo nó. Nhìn xa hơn vào vũ trụ nghĩa là nhìn xa hơn về quá khứ, nên các quan sát ở những khoảng cách khác nhau tái dựng phần sớm hơn của đường cong giãn nở chứ không chỉ đơn thuần giá trị của nó hôm nay.}
\end{classroomqa}

Bài giảng cho biết rằng các quan sát thời muộn sẵn có vào lúc đó phù hợp với một nhánh mũ ở mức xấp xỉ 5 phần trăm và mô tả chứng cứ là thuyết phục, dù không chính xác đến mức của một phép đo chính xác như các tính chất của electron. Điều này nên được hiểu là đánh giá mang tính lịch sử của bài giảng, chứ không phải một phát biểu về độ chính xác ở thời điểm hiện nay.

\begin{classroomqa}
\audiencequestion{Liệu đường cong giãn nở uốn lên đó cuối cùng có ngăn chúng ta lần ngược các quan sát về tới Big Bang hay phông vi sóng vũ trụ không?\lecturetimestamp{00:55:05}}
\lecturerresponse{Điều đó không được theo đuổi như một câu hỏi về việc ngoại suy đồ thị. Vấn đề bất biến thích hợp là sự tồn tại của một chân trời sự kiện, và đó là câu hỏi nhân quả mà bây giờ ta chuyển sang.}
\end{classroomqa}

\section{Các hệ liên kết và tương lai xa}

Nếu sự giãn nở mũ tiếp tục, mọi thành phần bị pha loãng đều trở nên không đáng kể so với \(\rho_{\rm vac}\). Điều đó không có nghĩa là bản thân các nguyên tử, các hệ Mặt Trời hay các thiên hà bị kéo giãn tách rời nhau. Những hệ được giữ với nhau đủ mạnh để vẫn liên kết thì không đi theo dòng Hubble.

\begin{classroomqa}
\audiencequestion{Các thiên hà hiện nay không giãn nở rồi cuối cùng có giãn nở khi năng lượng chân không chi phối không?\lecturetimestamp{00:57:44}}
\lecturerresponse{Các hệ liên kết thì vẫn liên kết. Dải Ngân Hà và Andromeda thuộc về cấu trúc lân cận sẽ còn ở cùng nhau, trong khi những thiên hà đủ xa đã chuyển động theo sự giãn nở thì tiếp tục đi ra xa và cuối cùng vượt ra ngoài chân trời của chúng ta.}
\end{classroomqa}

Sự phân biệt giữa \emph{bị hút} và \emph{liên kết} là thiết yếu. Một tên lửa được phóng vượt quá vận tốc thoát vẫn còn ở dưới tác dụng hấp dẫn hút của Trái Đất, nhưng nó không quay lại. Tương tự, lực hấp dẫn thông thường có thể là lực cục bộ lớn nhất giữa hai thiên hà mà quỹ đạo tương đối của chúng vẫn không phải là quỹ đạo liên kết.

\begin{classroomqa}
\audiencequestion{Thiên hà của chúng ta có liên hệ hấp dẫn với một nhóm. Liệu sự giãn nở gia tốc cuối cùng có thắng được lực hút đó không?\lecturetimestamp{00:59:17}}
\lecturerresponse{Chỉ phần thật sự liên kết của cấu trúc lân cận mới còn ở cùng nhau. Một cụm chỉ như một nhãn mô tả thì chưa đủ; liên kết có nghĩa là chuyển động tương đối nằm dưới vận tốc thoát. Các thiên hà xa hơn có thể bị hút mà vẫn lùi xa với tốc độ vượt vận tốc thoát.}
\end{classroomqa}

Để trực giác hơn, bài giảng biểu diễn gia tốc xuyên tâm tương đối một cách sơ đồ như tổng của một hạng hút nghịch đảo bình phương và một hạng chân không tuyến tính theo độ tách:
\begin{equation}
\ddot r\sim-\frac{G_NM}{r^2}+H^2r.
\label{eq:l08-newtonian-vacuum}
\end{equation}
Các hệ số là độc lập với thời gian trong giới hạn de Sitter lý tưởng hóa. Hạng thứ hai chỉ lớn lên vì \(r\) tăng; nó không phát sinh từ sự suy yếu dần của \(G_N\).

\editorialnote{Phương trình~\eqref{eq:l08-newtonian-vacuum} là một mô tả Newton mang tính heuristic về gia tốc tương đối, không phải một định luật lực bổ sung cơ bản của thuyết tương đối rộng. Trong không gian de Sitter chính xác, các đường thế giới rơi tự do lân cận có gia tốc tương đối tỉ lệ với \(H^2r\).}

\begin{figure}[htbp]
\centering
\includegraphics[width=0.88\textwidth]{lecture_08_figure_09.png}
\caption{So sánh trên bảng giữa lực hút nghịch đảo bình phương và lực đẩy chân không hiệu dụng tăng tuyến tính theo độ tách.}
\lectureframe{01:04:10.000}
\end{figure}

\begin{classroomqa}
\audiencequestion{Bản thân hiệu ứng hướng ra ngoài có tăng theo thời gian không?\lecturetimestamp{01:01:18}}
\lecturerresponse{Không, nếu xét tại một độ tách cố định. Hệ số của nó là hằng số. Các vật thể đang lùi xa tự do thì ngày càng cách nhau xa hơn, và khi đó hạng tỉ lệ với \(r\) lớn hơn vì khoảng cách của chúng lớn hơn.}
\end{classroomqa}

\begin{classroomqa}
\audiencequestion{Có thể mô tả hiệu ứng tổng thể như một sự giảm dần của hằng số Newton không?\lecturetimestamp{01:02:30}}
\lecturerresponse{Không. \(G_N\) vẫn là hằng số. Ở một độ tách cho trước, cả hai hệ số đều không đổi; chỉ có tầm quan trọng tương đối của lực hút \(1/r^2\) và của hạng tỉ lệ với \(r\) thay đổi theo khoảng cách.}
\end{classroomqa}

\begin{classroomqa}
\audiencequestion{Có phải những thiên hà ấy đã liên kết hấp dẫn trong vũ trụ sơ khai, trước khi năng lượng chân không trở nên quan trọng không?\lecturetimestamp{01:05:03}}
\lecturerresponse{Sự chi phối của lực hút thông thường tự nó không hàm ý có liên kết. Nhiều thiên hà đã tách xa nhau với tốc độ xấp xỉ vận tốc thoát hoặc lớn hơn. Về sau năng lượng chân không quyết định xem sự giãn nở quy mô lớn chậm dần tới tái co sụp hay gia tốc mà đi xa, nhưng việc lực là lớn nhất và việc quỹ đạo là liên kết vẫn là hai khẳng định khác nhau.}
\end{classroomqa}

Vì thế, một người quan sát ở tương lai rất xa có thể sống trong một cấu trúc lân cận đã sáp nhập trong khi mọi thiên hà không liên kết đều đã trở nên không thể quan sát. Không có ghi chép lịch sử, người quan sát đó hoàn toàn có thể suy ra một vũ trụ gần như trống rỗng. Sự cô đơn nhân quả này, chứ không phải sự phá hủy các vật thể địa phương, là phần dẫn nhập vật lý cho chân trời.

\section{Ý nghĩa của chân trời sự kiện}

Sau giờ nghỉ, đặt
\begin{equation}
a(t)=e^{Ht}
\end{equation}
và viết metric thành
\begin{equation}
ds^2=-dt^2+e^{2Ht}d\mathbf{x}^2.
\label{eq:l08-de-sitter-metric}
\end{equation}

\begin{classroomqa}
\audiencequestion{Tên của không-thời gian được mô tả bởi metric này là gì?\lecturetimestamp{01:08:13}}
\lecturerresponse{Đó là không-thời gian de Sitter, thường được gọi ngắn gọn là không gian de Sitter. Các chân trời của nó quan trọng vì một vũ trụ bị năng lượng chân không chi phối tiến gần tới hình học này ở thời muộn.}
\end{classroomqa}

Trước khi vẽ cấu trúc nhân quả của nó, hãy nhắc lại vai trò của biểu đồ Penrose. Các khoảng vô hạn được nén vào một hình hữu hạn trong khi các tia null vẫn nằm ở \(45^\circ\). Trong không-thời gian phẳng, các vật thể có khối lượng kết thúc ở vô cùng kiểu thời tương lai, các tia null kết thúc ở vô cùng null, và một người quan sát quán tính được phép chờ lâu tùy ý thì không có chân trời sự kiện nào loại trừ vĩnh viễn một vùng không-thời gian.

\begin{figure}[htbp]
\centering
\includegraphics[width=0.80\textwidth]{lecture_08_figure_10.png}
\caption{Biểu đồ Penrose của không-thời gian phẳng dùng để phân biệt các đích kiểu thời và null trước khi đưa vào một chân trời.}
\lectureframe{01:10:35.000}
\end{figure}

Biểu đồ hố đen tạo ra sự tương phản. Một người quan sát ở bên ngoài chỉ có thể nhận tín hiệu từ các sự kiện nằm trong quá khứ nhân quả của phần tương lai trên đường thế giới của mình. Ngay cả khi thời gian quan sát tiến tới vô hạn, vùng bên trong vẫn bị loại trừ. Biên giữa những sự kiện có thể liên lạc với người quan sát đó và những sự kiện vĩnh viễn không thể là chân trời sự kiện.

\begin{definition}Biên của tập các sự kiện mà từ đó một tín hiệu nhân quả hướng về tương lai có thể, vào một lúc nào đó, tới được đường thế giới của một người quan sát chính là chân trời sự kiện tương lai của người quan sát ấy.
\label{def:l08-event-horizon}
\end{definition}

\begin{figure}[htbp]
\centering
\includegraphics[width=0.80\textwidth]{lecture_08_figure_11.png}
\caption{So sánh với hố đen: các nón ánh sáng lùi về quá khứ từ những quan sát bên ngoài ở các thời điểm muộn tùy ý không bao giờ bao gồm vùng ở phía bên kia chân trời sự kiện.}
\lectureframe{01:13:50.000}
\end{figure}

Vì vậy, chân trời sự kiện là một biên nhân quả toàn cục, không chỉ đơn thuần là khoảng cách mà tại đó hôm nay một vật tình cờ trông mờ nhạt. Muốn xác định nó, ta phải biết những tín hiệu nào có thể đến nơi ngay cả sau khi chờ lâu tùy ý.

\section{Nén tương lai vô hạn}

Các tọa độ không gian trong Phương trình~\eqref{eq:l08-de-sitter-metric} đã trải trên một lát cắt phẳng vô hạn. Ta sẽ giữ nguyên chúng và chỉ nén thời gian. Đưa vào thời gian conformal \(T\) bằng cách yêu cầu
\begin{equation}
dt^2=e^{2Ht}dT^2.
\label{eq:l08-conformal-definition}
\end{equation}
Lấy căn bậc hai theo hướng tương lai rồi sắp xếp lại, ta được
\begin{equation}
dT=e^{-Ht}dt.
\end{equation}

\begin{classroomqa}
\audiencequestion{Nguyên hàm đó là gì, kể cả dấu của nó?\lecturetimestamp{01:17:25}}
\lecturerresponse{Nó là \(-H^{-1}e^{-Ht}\). Một hằng số cộng chỉ dịch gốc của \(T\), nên hãy chọn nó sao cho tương lai vô hạn nằm tại \(T=0\).}
\end{classroomqa}

Do đó
\begin{equation}
T=-\frac{1}{H}e^{-Ht},
\qquad -\infty<T<0.
\label{eq:l08-conformal-time}
\end{equation}
Khi \(t\to-\infty\), \(T\to-\infty\). Khi thời gian riêng \(t\to+\infty\), \(T\to0^{-}\). Một tương lai vô hạn đã được đặt lên đường nằm ngang hữu hạn \(T=0\); không người quan sát nào đạt tới đó trong thời gian riêng hữu hạn.

\begin{figure}[htbp]
\centering
\includegraphics[width=0.86\textwidth]{lecture_08_figure_12.png}
\caption{Tích phân thời gian conformal và phép ánh xạ tương lai vô hạn \(t\to+\infty\) thành \(T=0^{-}\).}
\lectureframe{01:19:10.000}
\end{figure}

Từ Phương trình~\eqref{eq:l08-conformal-time},
\begin{equation}
HT=-e^{-Ht},
\qquad
e^{2Ht}=\frac{1}{H^2T^2}.
\label{eq:l08-conformal-factor}
\end{equation}
Dùng \(dt=e^{Ht}dT\) trong Phương trình~\eqref{eq:l08-de-sitter-metric} thì thu được
\begin{equation}
ds^2=\frac{1}{H^2T^2}
\left(-dT^2+d\mathbf{x}^2\right).
\label{eq:l08-conformal-metric}
\end{equation}
Biểu thức trong ngoặc là metric Minkowski. Toàn bộ hình học de Sitter giờ nằm trong một nhân tử conformal chung, nhân tử này phân kỳ khi \(T\to0^{-}\).

\begin{figure}[htbp]
\centering
\includegraphics[width=0.87\textwidth]{lecture_08_figure_13.png}
\caption{Phép biến đổi đại số liên hệ \(e^{2Ht}\) với \(1/(H^2T^2)\), chuẩn bị cho dạng conformal của metric.}
\lectureframe{01:22:50.000}
\end{figure}

\begin{classroomqa}
\audiencequestion{Điều kiện nào xác định quỹ đạo của một tia sáng trong metric này?\lecturetimestamp{01:23:49}}
\lecturerresponse{Một tia sáng đi theo một quỹ đạo null, \(ds^2=0\). Vì thế nhân tử conformal chung triệt tiêu, để lại \(dT=\pm dx\) dọc theo một phương không gian đã chọn.}
\end{classroomqa}

Các tia sáng là những đường thẳng \(45^\circ\) trong biểu đồ \((T,x)\):
\begin{equation}
\frac{dx}{dT}=\pm1.
\label{eq:l08-null-lines}
\end{equation}
Bức tranh tọa độ bị nén mạnh gần \(T=0\), nhưng nó vẫn bảo toàn chính xác điều ta cần: những sự kiện nào có thể liên lạc với nhau.

\begin{figure}[htbp]
\centering
\includegraphics[width=0.86\textwidth]{lecture_08_figure_14.png}
\caption{Mêtric bảo giác tại thời điểm người ta dùng điều kiện null \(ds^2=0\) để khôi phục các tia sáng \(45^\circ\).}
\lectureframe{01:24:05.000}
\end{figure}

\section{Alice và Bob trong không gian de Sitter}

Đặt Alice tại vị trí đồng chuyển động cố định \(x=0\). Đường thế giới của cô đi thẳng đứng lên trên và kết thúc trên biểu đồ tại \(T=0\), mặc dù đồng hồ đeo tay của cô cần thời gian riêng vô hạn mới tới được đó. Hãy vẽ các tia null hướng về quá khứ từ điểm cuối ấy. Chúng xác định miền lớn nhất mà từ đó Alice còn có thể nhận được một tín hiệu. Các sự kiện ở bên ngoài miền đó sẽ vĩnh viễn không thể tiếp cận.

\begin{figure}[htbp]
\centering
\includegraphics[width=0.88\textwidth]{lecture_08_figure_15.png}
\caption{Nón ánh sáng quá khứ giới hạn của Alice. Biên null của nó là chân trời sự kiện de Sitter phụ thuộc người quan sát của cô.}
\lectureframe{01:28:50.000}
\end{figure}

Chân trời này là riêng của từng người quan sát. Một người quan sát khác có một điểm cuối tương lai khác và vì thế có một nón ánh sáng quá khứ giới hạn khác. Giả sử có lúc Bob nằm trong miền nhân quả của Alice và hai người trao đổi thông điệp. Nếu về sau Bob đi theo dòng Hubble băng qua chân trời của Alice, thì các tín hiệu phát ra sau lần băng qua đó sẽ không thể tới được Alice, bất kể cô chờ bao lâu. Bob có chân trời riêng của mình và cuối cùng cũng mất đi khả năng tương ứng là nhận các tín hiệu phát muộn từ Alice.

\begin{figure}[htbp]
\centering
\includegraphics[width=0.88\textwidth]{lecture_08_figure_16.png}
\caption{Alice và Bob với các biên nhân quả riêng biệt của họ. Một Bob đồng chuyển động có thể rời khỏi quá khứ nhân quả của Alice mặc dù trước đó họ đã trao đổi tín hiệu.}
\lectureframe{01:30:20.000}
\end{figure}

Nếu Alice và Bob nắm tay nhau, hoặc thực tế hơn là vẫn bị liên kết ở khoảng cách riêng cố định, thì độ tách trong tọa độ đồng chuyển động của họ co lại như \(a^{-1}\). Các đường cong của họ tiến lại gần nhau trên biểu đồ bảo giác, nhưng điều này không có nghĩa là khoảng cách riêng của họ biến mất. Hệ số bảo giác tăng dần bù cho khe hở tọa độ đang co lại, và việc liên lạc vẫn có thể duy trì vô thời hạn.

\begin{classroomqa}
\audiencequestion{Sau khi các chân trời tách họ ra, Bob có thể nhìn thấy Alice vào cùng một thời điểm không?\lecturetimestamp{01:31:43}}
\lecturerresponse{Không người quan sát nào từng nhìn thấy một sự kiện ở xa "đồng thời" theo nghĩa đen đó; quan sát luôn là quan sát một sự kiện nằm trong nón ánh sáng quá khứ. Phát biểu sắc nét hơn là sau khi băng qua chân trời, không tín hiệu mới phát nào từ đường thế giới kia có thể tới nơi.}
\end{classroomqa}

Trong ứng dụng vũ trụ học ở tương lai xa, Alice biểu diễn tàn dư địa phương bị liên kết của chúng ta còn Bob biểu diễn từng thiên hà xa không bị liên kết. Các thiên hà không đột ngột vụt tắt tại một thời điểm hữu hạn trên đồng hồ của Alice. Những tín hiệu về sau của chúng bị kéo giãn trên các khoảng thời gian dài hơn và dịch sang các bước sóng dài hơn cho đến khi chúng trở nên không thể phát hiện được về mặt vật lý.

\begin{classroomqa}
\audiencequestion{Lý thuyết thực sự có nói rằng sẽ không còn thiên hà xa nào nhìn thấy được trong tương lai xa không?\lecturetimestamp{01:33:38}}
\lecturerresponse{Các thiên hà không bị liên kết đi ra ngoài chân trời sự kiện của người quan sát địa phương. Những tín hiệu nhìn thấy được cuối cùng của chúng đến với độ trễ ngày càng lớn và bị dịch đỏ ngày càng mạnh, chỉ còn lại cấu trúc địa phương bị liên kết là còn quan sát được cho các mục đích thực tế.}
\end{classroomqa}

\section{Độ dịch đỏ tại chân trời}

Người ta vẫn có thể vẽ một tín hiệu ánh sáng cũ từ Bob tới Alice tại một thời bảo giác rất muộn, điều này có thể tạo cảm giác rằng cô vẫn tiếp tục nhìn thấy anh. Nội dung vật lý nằm ở khoảng cách thời gian và năng lượng của các tín hiệu ấy. Những khoảng thời gian bằng nhau trên đồng hồ của Bob tới nơi với khoảng cách theo thời gian riêng của Alice ngày càng lớn hơn. Tần số và năng lượng photon giảm:
\begin{equation}
\nu_{\rm received}\longrightarrow0,
\qquad
E_{\rm received}=h\nu_{\rm received}\longrightarrow0.
\end{equation}
Ảnh của Bob chậm dần, đỏ hơn, rồi mờ dần. Tín hiệu giới hạn mang lượng thông tin hữu dụng triệt tiêu.

\begin{figure}[htbp]
\centering
\includegraphics[width=0.89\textwidth]{lecture_08_figure_17.png}
\caption{Trạng thái cuối trên bảng so sánh độ dịch đỏ ở chân trời de Sitter với cách nhìn từ bên ngoài quen thuộc về chân trời lỗ đen.}
\lectureframe{01:36:20.000}
\end{figure}

\begin{classroomqa}
\audiencequestion{Điều này có giống lỗ đen, nơi một người quan sát bên ngoài không bao giờ thấy một vật đang rơi vào băng qua chân trời không?\lecturetimestamp{01:36:09}}
\lecturerresponse{Có. Trong cả hai trường hợp, người quan sát nhận được các tín hiệu ngày càng bị dịch đỏ và không bao giờ trực tiếp thấy việc băng qua hoàn tất. Biểu đồ nhân quả cho biết những phát xạ về sau nào sẽ không bao giờ tới nơi; còn độ dịch đỏ giải thích thông tin quan sát được cuối cùng mờ dần như thế nào.}
\end{classroomqa}

Phép loại suy này gợi ra một câu hỏi về nhiệt độ.

\begin{classroomqa}
\audiencequestion{Làm thế nào Alice có thể đo nhiệt độ ngay phía trên một chân trời?\lecturetimestamp{01:39:34}}
\lecturerresponse{Về mặt thao tác, cô có thể hạ một nhiệt kế xuống bằng dây cáp, giữ nó không cho băng qua, rồi kéo nó lên. Một người quan sát de Sitter có thể hình dung thí nghiệm tương tự gần chân trời vũ trụ học. Trong lý thuyết lượng tử, các chân trời là nóng, nhưng bài giảng hiện tại dừng ở hình học cổ điển và hoãn phép tính đó lại.}
\end{classroomqa}

Ở đây không suy ra công thức nào cho nhiệt độ de Sitter. Bức xạ Hawking chỉ được nêu ra như một báo trước về hiệu chỉnh lượng tử đối với bức tranh cổ điển này.

\section{Bán kính chân trời cố định}

Biểu đồ bảo giác có thể được chuyển trở lại tọa độ thời gian ban đầu. Hai điểm đồng chuyển động ở độ tách tọa độ cố định thỏa mãn
\begin{equation}
D(t)=e^{Ht}\lvert\Delta x\rvert,
\end{equation}
nên khoảng cách riêng của chúng tăng theo hàm mũ.

\begin{classroomqa}
\audiencequestion{Cùng chuyển động đó trông như thế nào trong tọa độ thời gian ban đầu, chưa biến đổi?\lecturetimestamp{01:41:14}}
\lecturerresponse{Với các tọa độ đồng chuyển động cố định, toàn bộ mêtric không gian tăng theo \(e^{2Ht}\); do đó các độ tách riêng tăng như \(e^{Ht}\). Việc đảo phép biến đổi tọa độ làm cho sự giãn nở ấy hiện ra tường minh.}
\end{classroomqa}

Chân trời ứng xử khác đi vì khoảng cách tọa độ đồng chuyển động từ nó tới Alice không cố định. Từ đường null \(45^\circ\) kết thúc tại \(T=0\), khoảng cách tới chân trời trên một lát \(T<0\) không đổi là
\begin{equation}
\Delta x_H=-T=\frac{1}{H}e^{-Ht}.
\label{eq:l08-coordinate-horizon}
\end{equation}
Nhân với thừa số tỉ lệ, ta được
\begin{equation}
D_H(t)=a(t)\Delta x_H
=e^{Ht}\frac{e^{-Ht}}{H}
=\frac{1}{H}.
\label{eq:l08-fixed-horizon}
\end{equation}
Khôi phục đơn vị, \(D_H=c/H\). Khe hở tọa độ co lại, thang của mêtric tăng lên, và tích của chúng vẫn không đổi.

\begin{classroomqa}
\audiencequestion{Khi thời gian trôi qua, khoảng cách từ Alice tới chân trời tăng, giảm hay giữ nguyên?\lecturetimestamp{01:42:42}}
\lecturerresponse{Nó giữ nguyên. Khoảng cách tọa độ của nó co lại trong khi hệ số bảo giác tăng lên, để lại bán kính chân trời theo khoảng cách riêng cố định ở \(H^{-1}\) trong hệ đơn vị với \(c=1\).}
\end{classroomqa}

\begin{figure}[H]
\centering
\includegraphics[width=0.90\textwidth]{lecture_08_figure_18.png}
\caption{Trạng thái cuối trên bảng: các lát ngang kế tiếp có khoảng cách tọa độ tới chân trời co nhỏ dần nhưng cùng một khoảng cách riêng.}
\lectureframe{01:43:20.000}
\end{figure}

\begin{classroomqa}
\audiencequestion{Alice có tiếp tục nhìn thấy quá khứ được chiếu lên chân trời vì những ảnh cũ không bao giờ biến mất hẳn không?\lecturetimestamp{01:43:45}}
\lecturerresponse{Về mặt cổ điển, các ảnh chỉ tồn tại theo nghĩa giới hạn là chúng ngày càng khó phát hiện hơn do độ dịch đỏ không bị chặn. Phát biểu đó bỏ qua bức xạ Hawking gắn với chân trời, vốn sẽ làm thay đổi mô tả lượng tử.}
\end{classroomqa}

Vì vậy bài giảng khép lại ở một ranh giới xác định rõ ràng. Năng lượng chân không đã tạo ra sự giãn nở theo hàm mũ; sự giãn nở theo hàm mũ đã tạo ra một chân trời sự kiện cố định, phụ thuộc người quan sát; và các tín hiệu gần chân trời đó đã bị đẩy tới tần số triệt tiêu. Nhiệt độ của chân trời và bức xạ Hawking thuộc về phần tiếp theo. Giả sử, như Susskind nói, rằng chúng ta không băng qua chân trời của nhau, hẹn gặp lại ở bài giảng sau.
